\documentclass[12pt]{article}
\usepackage{amsmath, amssymb, amsthm}
\usepackage[utf8]{inputenc}
\usepackage[T1]{fontenc}
\usepackage{hyperref}
\usepackage{geometry}
\usepackage{CJKutf8}   % for Chinese

\geometry{margin=1in}

\newtheorem{theorem}{Theorem}
\newtheorem{lemma}[theorem]{Lemma}
\newtheorem{claim}[theorem]{Claim}
\newtheorem{proposition}[theorem]{Proposition}
\newtheorem{corollary}[theorem]{Corollary}
\theoremstyle{definition}
\newtheorem{definition}[theorem]{Definition}
\newtheorem{example}[theorem]{Example}
\theoremstyle{remark}
\newtheorem{remark}[theorem]{Remark}

\title{Plan: $\varepsilon$-Light Subsets in Graphs}
\author{}
\date{Last updated: \today}

\begin{document}
\maketitle

\begin{abstract}
We study the following question: does there exist a universal constant $c > 0$ such that for every graph $G = (V,E)$ and every $\varepsilon \in [0,1]$, the vertex set $V$ contains an $\varepsilon$-light subset $S$ of size at least $c\varepsilon|V|$?  We argue the answer is \textbf{yes}, sketch a proof decomposed into a sparse and a dense case, and identify the tight example that forces $c \le 1/2$.
\end{abstract}

% ----------------------------------------------------------------
\section{Problem Statement}

\begin{CJK}{UTF8}{gbsn}
\textbf{中文问题陈述.} 设 $G = (V,E)$ 为图，$n = |V|$。对于子集 $S \subseteq V$，定义 $G_S = (V, E(S,S))$ 为保留所有顶点但仅保留 $S$ 内部边的子图。设 $L$ 为 $G$ 的拉普拉斯矩阵，$L_S$ 为 $G_S$ 的拉普拉斯矩阵。若矩阵 $\varepsilon L - L_S$ 半正定（即 $L_S \preceq \varepsilon L$），则称 $S$ 为 \textbf{$\varepsilon$-轻集}。问：是否存在绝对常数 $c > 0$，使得对任意图 $G$ 和任意 $\varepsilon \in [0,1]$，存在大小至少为 $c\varepsilon|V|$ 的 $\varepsilon$-轻集？
\end{CJK}

\bigskip

\textbf{English statement.}
Let $G = (V,E)$ be a graph with $n = |V|$ vertices.  For $S \subseteq V$, let
$G_S = (V, E(S,S))$ be the graph on the same vertex set but with only the edges
whose \emph{both} endpoints lie in $S$.  Let $L$ be the Laplacian of $G$ and $L_S$ be the Laplacian of $G_S$.

\begin{definition}[$\varepsilon$-light set]
A subset $S \subseteq V$ is \emph{$\varepsilon$-light} if
$L_S \preceq \varepsilon L$, i.e., the matrix $\varepsilon L - L_S$ is positive semidefinite.
\end{definition}

\textbf{Question.} Does there exist a constant $c > 0$ such that for every graph $G$ and every $\varepsilon \in [0,1]$, there exists an $\varepsilon$-light subset $S \subseteq V$ with $|S| \ge c\,\varepsilon\,|V|$?

% ----------------------------------------------------------------
\section{Spectral Reformulation}

The condition $L_S \preceq \varepsilon L$ is equivalent to
\[
  \forall x \in \mathbb{R}^n:\quad
  \sum_{(u,v)\in E(S,S)} (x_u - x_v)^2
  \;\le\;
  \varepsilon \sum_{(u,v)\in E} (x_u - x_v)^2.
\]
Informally, the Dirichlet energy of any signal $x$ on the induced subgraph $G[S]$ is at most an $\varepsilon$-fraction of its energy on all of $G$.

\textbf{Trivial observations.}
\begin{itemize}
  \item Any independent set $I$ satisfies $L_I = 0 \preceq \varepsilon L$ for all $\varepsilon \ge 0$.
  \item $S = V$ satisfies $L_V = L \preceq 1\cdot L$, so $V$ is $1$-light.
  \item Any single vertex $\{v\}$ is $\varepsilon$-light for all $\varepsilon$ (no intra-$S$ edges).
\end{itemize}

% ----------------------------------------------------------------
\section{Key Examples}

\subsection{Complete graph \texorpdfstring{$K_n$}{Kn}}

$L_{K_n} = nI - \mathbf{1}\mathbf{1}^\top$.
For $S \subseteq V$ with $|S| = k$, the graph $G_S$ is a copy of $K_k$ on vertices $S$ plus isolated vertices.  The maximum eigenvalue of $L_S$ equals $k$, while the maximum eigenvalue of $\varepsilon L$ equals $\varepsilon n$.  A careful check (using the eigenvectors of $L_S$ in the ambient $n$-dimensional space) shows:
\[
  L_S \preceq \varepsilon L \iff k \le \varepsilon n.
\]
Hence \emph{any} set of size $\lfloor\varepsilon n\rfloor$ is $\varepsilon$-light for $K_n$, achieving $c = 1$.

\subsection{Complete bipartite graph \texorpdfstring{$K_{n/2,\,n/2}$}{Knn}}

\label{sec:bipartite}

Let $A, B$ be the two sides, $|A|=|B|=n/2$.

\begin{claim}
For $\varepsilon \in [0,1)$, every $\varepsilon$-light set $S \subseteq V$ satisfies $S \subseteq A$ or $S \subseteq B$ (i.e.\ $S$ is ``pure'').
\end{claim}

\begin{proof}
Suppose $S$ contains vertices from both sides: let $a' = |S \cap A| \ge 1$ and $b' = |S \cap B| \ge 1$.
Take $x$ supported only on $B' = S \cap B$ (i.e.\ $x_i = 0$ for $i \notin B'$).
Then
\[
  x^\top L_S x = \sum_{\substack{i \in S \cap A,\; j \in B'}} (x_i - x_j)^2
               = a' \sum_{j \in B'} x_j^2.
\]
\[
  x^\top L\, x = \sum_{i \in A,\; j \in B} (x_i - x_j)^2
               = \frac{n}{2} \sum_{j \in B'} x_j^2.
\]
The ratio is $x^\top L_S x \,/\, x^\top L\,x = a'/(n/2)$.
For $L_S \preceq \varepsilon L$ we need $a'/(n/2) \le \varepsilon$, i.e.\ $a' \le \varepsilon n/2$.

By the symmetric argument (taking $x$ supported on $A' = S \cap A$) we also need $b' \le \varepsilon n/2$.

Consider the bipartite unit vector: $x_i = 1/\sqrt{a'+b'}$ for $i \in A \cap S$,
$x_j = -1/\sqrt{a'+b'}$ for $j \in B'$, zero elsewhere.  Then
\begin{align*}
  x^\top L_S x &= \frac{4\,a'b'}{a'+b'},\\
  x^\top L\,x &= \frac{2\,a'b'}{a'+b'} + \frac{n}{2}.
\end{align*}
For $L_S \preceq \varepsilon L$:
\[
  \frac{4\,a'b'}{a'+b'} \le \varepsilon\!\left(\frac{2\,a'b'}{a'+b'} + \frac{n}{2}\right)
  \implies
  a'+b' \le \frac{\varepsilon n}{2-\varepsilon}.
\]

Taking $a' \to n/2$ (the full side $A$) gives $b' \le \frac{\varepsilon n}{2-\varepsilon} - \frac{n}{2} = \frac{(\varepsilon-1)n/2}{2-\varepsilon}$, which is \emph{negative} for $\varepsilon < 1$.  Hence no mixed set with $a' = n/2$ is $\varepsilon$-light when $\varepsilon < 1$, and by a continuity argument the same holds for any $a' \ge 1, b' \ge 1$ with $a' + b' \le n/2$ only if $a' = b' = 0$... In fact the necessary condition $b' \le (\varepsilon - 1)n/2 < 0$ (for $\varepsilon < 1$) forces $b' = 0$, a contradiction.
\end{proof}

\begin{corollary}
For $K_{n/2,n/2}$ and $\varepsilon \in [0,1)$, the maximum size of an $\varepsilon$-light set is $n/2$ (take $S = A$ or $S = B$, giving $L_S = 0$).
\end{corollary}

\textbf{Tightness for $c$.}
For $\varepsilon \in [0,1)$ the maximum $\varepsilon$-light set has size $n/2$.  We need $n/2 \ge c\varepsilon n$, i.e.\ $c \le 1/(2\varepsilon)$.  As $\varepsilon \to 1^-$ this forces $c \le 1/2$.

\begin{proposition}
Any valid constant $c$ must satisfy $c \le 1/2$.
\end{proposition}

\subsection{Star \texorpdfstring{$K_{1,n-1}$}{K1n-1}}

Taking $S = V \setminus \{\text{hub}\}$ yields $L_S = 0$ and $|S| = n-1 \approx n \ge c\varepsilon n$.

\subsection{Path \texorpdfstring{$P_n$}{Pn}}

The path is very sparse ($m = n-1$).  Any $S$ has at most $|S|-1$ intra-$S$ edges, so $L_S \preceq (|S|-1)/(n-1) \cdot L$ (rough bound).  Taking $|S| = 1 + \varepsilon(n-1) \approx \varepsilon n$ works.

\subsection{Expander graphs}

A $d$-regular Ramanujan graph has spectral gap $\lambda_2 \ge d - 2\sqrt{d-1}$.  By the expander mixing lemma, for any $S$ with $|S| = k$:
$e(S,S) \le \frac{dk^2}{2n} + \frac{\lambda}{2}\cdot k$, where $\lambda = d - \lambda_2$.
This controls the number of intra-$S$ edges; combined with spectral domination arguments, one can show that $S$ of size $\varepsilon n$ is approximately $\varepsilon$-light (up to spectral gap corrections).

% ----------------------------------------------------------------
\section{Main Result}

\begin{theorem}[Main claim]\label{thm:main}
There exists an absolute constant $c > 0$ (and $c = 1/2$ is believed to be achievable, and tight due to $K_{n/2,n/2}$) such that for every graph $G = (V,E)$ and every $\varepsilon \in [0,1]$, there exists an $\varepsilon$-light subset $S \subseteq V$ with
\[
  |S| \ge c\,\varepsilon\,|V|.
\]
\end{theorem}

The proof splits into a sparse case and a dense case.

% ----------------------------------------------------------------
\section{Proof Strategy}

Let $n = |V|$, $m = |E|$, and $d_{\mathrm{avg}} = 2m/n$.

\subsection{Sparse case: \texorpdfstring{$d_{\mathrm{avg}} \le 1/(2\varepsilon)$}{sparse}}

\begin{lemma}[Greedy independent set]
Every graph on $n$ vertices with average degree $d$ has an independent set of size at least $n/(d+1)$.
\end{lemma}

\begin{proof}
Process vertices in any order; greedily include a vertex if none of its already-included neighbors is selected.  At most $d+1$ vertices are "blocked" per included vertex.
\end{proof}

\begin{corollary}
In the sparse case, the greedy independent set $I$ satisfies
\[
  |I| \;\ge\; \frac{n}{d_{\mathrm{avg}}+1} \;\ge\; \frac{n}{\frac{1}{2\varepsilon}+1} \;\ge\; \frac{\varepsilon n}{1 + 2\varepsilon} \;\ge\; \frac{\varepsilon n}{3}.
\]
Since $L_I = 0 \preceq \varepsilon L$, the set $I$ is $\varepsilon$-light.  Taking $c = 1/3$ suffices here.
\end{corollary}

\subsection{Dense case: \texorpdfstring{$d_{\mathrm{avg}} > 1/(2\varepsilon)$}{dense}}

\paragraph{Construction.}
Include each vertex in $S$ independently with probability $p = \sqrt{\varepsilon}$.

\paragraph{Expected size.}
$\mathbb{E}[|S|] = \sqrt{\varepsilon}\cdot n$.  Since $\sqrt{\varepsilon} \ge \varepsilon$ for $\varepsilon \in [0,1]$,
we have $\mathbb{E}[|S|] \ge \varepsilon n$.

\paragraph{Expected Laplacian.}
Each edge $(u,v)$ appears in $E(S,S)$ iff both $u,v \in S$, which happens with probability $p^2 = \varepsilon$.  Therefore
\[
  \mathbb{E}[L_S] \;=\; p^2 L \;=\; \varepsilon L.
\]

\paragraph{Concentration: size.}
Let $X = |S| = \sum_v \mathbf{1}[v \in S]$.  By the Paley--Zygmund inequality,
\[
  \Pr\!\left(X \ge \tfrac{1}{2}\mathbb{E}[X]\right)
  \;\ge\; \frac{(\mathbb{E}[X])^2}{4\,\mathbb{E}[X^2]}.
\]
One computes $\mathbb{E}[X^2] = pn + p^2 n(n-1)$.  For large $n$ and $p = \sqrt{\varepsilon}$:
\[
  \Pr\!\left(|S| \ge \tfrac{\sqrt{\varepsilon}}{2}\,n\right) \;\ge\; \frac{1}{4}\cdot \frac{1}{1 + 1/(p(n-1))} \;\to\; \frac{1}{4}\quad\text{as }n\to\infty.
\]
Since $\sqrt{\varepsilon}/2 \ge \varepsilon/2$, with probability $\ge 1/4$ we have $|S| \ge \varepsilon n/2$.

\paragraph{Concentration: PSD condition.}
We need to show $L_S \preceq \varepsilon L$ (or $L_S \preceq 2\varepsilon L$) with positive probability.

Write $L_S = \sum_{(u,v)\in E} X_u X_v L_{uv}$, where $X_v \overset{\mathrm{iid}}{\sim} \mathrm{Bern}(p)$ and $L_{uv} = (e_u - e_v)(e_u - e_v)^\top$.

\textbf{Martingale approach.}  Expose $X_1, \ldots, X_n$ one by one.  Let $\mathcal{F}_k = \sigma(X_1,\ldots,X_k)$ and
\[
  M_k = \mathbb{E}[L_S \mid \mathcal{F}_k].
\]
This is a matrix-valued martingale.  The step differences satisfy
\[
  \|M_k - M_{k-1}\|_{\mathrm{op}}
  \;\le\; (X_k - p)\sum_{\substack{v:\,(v,k)\in E}} L_{kv}
  \;\preceq\; d_k \|L\|_{\mathrm{op}},
\]
where $d_k = \deg(k)$.

By the \textbf{Matrix Azuma inequality} (see Tropp 2012, \emph{User-Friendly Tail Bounds for Sums of Random Matrices}), for any $t > 0$:
\[
  \Pr\!\left(\lambda_{\max}\!\left(L_S - \varepsilon L\right) \ge t\right)
  \;\le\; n \cdot \exp\!\left(-\frac{t^2/2}{\sigma^2}\right),
\]
where $\sigma^2 = \sum_k d_k^2 \|L\|_{\mathrm{op}}^2$.  In the \emph{dense} case $d_{\mathrm{avg}} \ge 1/(2\varepsilon)$, the operator norm $\|L\|_{\mathrm{op}} = \lambda_{\max}(L) \ge d_{\mathrm{avg}} \ge 1/(2\varepsilon)$, and the bound $t = \varepsilon \lambda_{\max}(L)$ is large relative to $\sigma$, giving a useful tail bound.

\textbf{Combining.}  Let $E_1 = \{|S| \ge \varepsilon n/2\}$ and $E_2 = \{L_S \preceq 2\varepsilon L\}$.  We have:
\begin{align*}
  \Pr(E_1) &\ge 1/4,\\
  \Pr(E_2^c) &\le n \cdot \exp(-\Omega(\varepsilon^2 n / d_{\max}^2)).
\end{align*}
For $\varepsilon$ and $n$ large enough relative to $d_{\max}^2$, $\Pr(E_2) \ge 3/4$, so
\[
  \Pr(E_1 \cap E_2) \ge \Pr(E_1) + \Pr(E_2) - 1 \ge 1/4 + 3/4 - 1 = 0.
\]
In fact the bound is strictly positive, so there exists a realization of $S$ satisfying both $|S| \ge \varepsilon n/2$ and $L_S \preceq 2\varepsilon L$.  (For a $c\varepsilon n$ bound with $c = 1/4$, replace $\varepsilon$ by $\varepsilon/2$ throughout.)

\begin{remark}
The argument above is not tight.  A more refined analysis (using decoupling inequalities for U-statistics, or a different $p$) should recover $c = 1/2$.
\end{remark}

% ----------------------------------------------------------------
\section{Unifying the Two Cases}

\begin{theorem}[Sufficient condition for $c = 1/4$]\label{thm:c14}
For every graph $G$ and every $\varepsilon \in [0,1]$, there exists an $\varepsilon$-light subset $S$ with $|S| \ge (\varepsilon/4)\,|V|$.
\end{theorem}

\begin{proof}[Proof sketch]
\begin{itemize}
\item If $d_{\mathrm{avg}} \le 1/(2\varepsilon)$: apply the greedy independent set lemma to get $|I| \ge \varepsilon n / 3$.
\item If $d_{\mathrm{avg}} > 1/(2\varepsilon)$: apply the random construction with $p = \sqrt{\varepsilon}$, use the matrix concentration argument above, and take $c = 1/4$.
\end{itemize}
In both cases $|S| \ge (\varepsilon/4) n$.
\end{proof}

% ----------------------------------------------------------------
\section{Tightness: \texorpdfstring{$c \le 1/2$}{c <= 1/2}}

From Section~\ref{sec:bipartite}: for $K_{n/2,n/2}$ and $\varepsilon \in (0,1)$, the maximum $\varepsilon$-light set has size $n/2$.  The bound $|S| \ge c\varepsilon n$ requires $c \le n/2 / (\varepsilon n) = 1/(2\varepsilon)$, which as $\varepsilon \to 1^-$ gives $c \le 1/2$.

\textbf{Conjecture.} The tight constant is $c = 1/2$, i.e.\ for every graph $G$ and every $\varepsilon \in [0,1]$, there exists an $\varepsilon$-light $S$ with $|S| \ge \varepsilon n/2$.

% ----------------------------------------------------------------
\section{Key Lemmas to be Made Rigorous}

\begin{enumerate}
  \item \textbf{Greedy independent set}: stated and proved above.
  \item \textbf{Paley--Zygmund for $|S|$}: standard second-moment calculation.
  \item \textbf{Matrix Azuma inequality} (Tropp 2012): for a matrix-valued martingale with bounded increments.
  \item \textbf{Decoupling for U-statistics} (de la Pe\~{n}a--Gin\'{e}): to improve the dependence in the dense-case concentration from $d_{\max}$ to $d_{\mathrm{avg}}$.
  \item \textbf{Optimal $p$}: find the value of $p$ (as a function of $\varepsilon$ and $G$) that maximizes $|S|$ while keeping $L_S \preceq \varepsilon L$.
\end{enumerate}

% ----------------------------------------------------------------
\section{Open Questions}

\begin{enumerate}
  \item \textbf{Tight constant}: Is $c = 1/2$ achievable for all graphs?
  \item \textbf{Tight example}: Is $K_{n/2,n/2}$ the unique extremal example?  Are there other graphs where the maximum $\varepsilon$-light set has size exactly $\varepsilon n/2$?
  \item \textbf{Algorithmic}: Can the set $S$ be found in polynomial time?
  \item \textbf{Weighted graphs}: Does the result extend to weighted graphs (where $L = \sum_{e} w_e L_e$)?
  \item \textbf{Directed graphs}: What is the right analogue for directed Laplacians?
\end{enumerate}

% ----------------------------------------------------------------
\section*{References}

\begin{itemize}
  \item J.~A.~Tropp, ``User-Friendly Tail Bounds for Sums of Random Matrices,'' \emph{Foundations of Computational Mathematics}, 12(4):389--434, 2012.
  \item V.~H.~de la Pe\~{n}a and E.~Gin\'{e}, \emph{Decoupling: From Dependence to Independence}, Springer, 1999.
  \item D.~Spielman and N.~Srivastava, ``Graph Sparsification by Effective Resistances,'' \emph{SIAM Journal on Computing}, 40(6):1913--1926, 2011.
\end{itemize}

\end{document}
